\begin{alist}
\item Επειδή $P(1) = 1^4 - 1^3 - 5 \cdot 1^2 + 7 \cdot 1 - 2 = 1 - 1 - 5 + 7 - 2 = 0$, ο αριθμός 1 είναι ρίζα του πολυωνύμου.

\item Οι πιθανές ακέραιες λύσεις του πολυωνύμου είναι οι διαιρέτες του σταθερού όρου του δηλαδή οι διαιρέτες του 2 που είναι οι αριθμοί $1, 2, -1, -2$.
Ο αριθμός 1 είδαμε ότι είναι ρίζα του πολυωνύμου. Εξετάζουμε αν κάποιος από τους άλλους αριθμούς είναι ρίζα. Είναι:
\begin{itemize}
\item $P(2) = 16 - 8 - 20 + 14 - 2 = 0$
\item $P(-1) = 1 + 1 - 5 - 7 - 2 = -12 \ne 0$
\item $P(-2) = 16 + 8 - 20 - 14 - 2 = -12 \ne 0$
\end{itemize}
Επομένως οι ακέραιες ρίζες του πολυωνύμου είναι οι αριθμοί 1 και 2.
\end{alist}
